\section{Genuine tangent partial derivatives of the whole original homogeneous initial ideal}

For an unchanged original holomorphic germ ideal $I$, let
$K=\operatorname{map}(\operatorname{originalTaylorAlgHom})(I)$ be
its entire actual formal extension and let
$J=\operatorname{actualHomogeneousInitialIdeal}(K)$. The latter is
the existing ideal generated by the literal lowest homogeneous forms
of \emph{all} nonzero formal members of $K$, not a selected leading
term or generator replacement. Polynomial membership below always
uses the genuine polynomial-to-series inclusion into this same $J$.

\begin{theorem}[True homogeneous components and literal differentiation retain the genuine initial ideal dictionary]
\label{hir:actual-homogeneous-initial-component-derivative-dictionaries}
\uses{hir:actual-homogeneous-initial-form-coefficients,hir:actual-original-formal-variable-power-extension}
\lean{Hironaka.HironakaActualHomogeneousInitialTangentDerivatives.homogeneousComponent_mem_actual_initial_of_mem_power,Hironaka.HironakaActualHomogeneousInitialTangentDerivatives.pderiv_homogeneousComponent_succ,Hironaka.HironakaActualHomogeneousInitialTangentDerivatives.pderiv_homogeneousComponent_zero}
\leanok
For every entire formal ideal $K$, if $A\in K$ vanishes in degrees
below $q$, its genuine degree-$q$ component belongs to the same
$\operatorname{actualHomogeneousInitialIdeal}(K)$. Actual formal
coordinate differentiation satisfies
\[
 \partial_i\operatorname{HC}_{q+1}(F)
    =\operatorname{HC}_q(\partial_i F),\qquad
 \partial_i\operatorname{HC}_0(F)=0.
\]
These are identities for the unchanged series and actual coefficient
maps, not an assumed initial-form derivative model.
\end{theorem}
\begin{proof}
For the membership assertion, zero series give zero components. For
nonzero $A$, its actual least coefficient degree is at least $q$,
because every lower coefficient vanishes by the genuine formal
variable-power kernel. If that degree equals $q$, the component is
exactly the existing literal initial form of this actual $A\in K$,
hence is one of the genuine span generators. If the degree is
larger, the component is zero by the actual order/component theorem.

For differentiation, the actual coefficient at exponent $\alpha$ of
$\partial_i F$ is the original coefficient at
$\alpha+e_i$ times $\alpha_i+1$. The total degree of that exponent is
$|\alpha|+1$. Thus selecting degree $q+1$ before differentiation is
exactly selecting degree $q$ afterward, coefficient by coefficient.
No exponent can have degree $-1$, so differentiating the actual
constant component gives zero. Extensional equality of the genuine
formal coefficient arrays proves both identities.
\end{proof}

\begin{theorem}[Original first-symbol closure transfers to the entire same homogeneous initial ideal and literal polynomial partial derivatives]
\label{hir:actual-original-first-symbol-whole-initial-pderiv-transfer}
\uses{hir:actual-analytic-ideal-member-finite-jet-lift,hir:actual-homogeneous-initial-component-derivative-dictionaries,hir:actual-whole-original-taylor-coordinate-derivative,hir:relative-normal-basis-actual-derivation-evaluation}
\lean{Hironaka.HironakaActualHomogeneousInitialTangentDerivatives.actual_homogeneous_initial_ideal_pderiv_of_first_symbol_closure,Hironaka.HironakaActualHomogeneousInitialTangentDerivatives.actual_polynomial_pderiv_mem_initial_of_first_symbol_closure}
\leanok
Fix an actual original coordinate $i$ and assume the genuine original
first-symbol membership criterion
$D_ig\in I+\mathfrak m^q$ for every positive $q$ and every whole
$g\in I\cap\mathfrak m^q$. Then every $H\in J$ satisfies
$\partial_i H\in J$. Consequently, for every actual original
polynomial $P$ whose genuine series inclusion belongs to this same
$J$, the inclusion of the literal polynomial
$\operatorname{MvPolynomial.pderiv}_i(P)$ belongs to the same $J$.
The criterion is derived from original geometry in the next theorem;
no formal derivative-stability certificate is substituted for it.
\end{theorem}
\begin{proof}
First treat every genuine initial generator $\operatorname{in}(F)$,
where $F$ is an arbitrary nonzero member of the entire same $K$.
If its actual least degree is zero, its derivative is zero. Otherwise
write that degree as $q=k+1$. The genuine finite approximation theorem
supplies $g\in I\cap\mathfrak m^q$ with the same entire degree-$q$
component as $F$. Original first-symbol closure writes
$D_ig=a+b$ with $a\in I$ and $b\in\mathfrak m^q$. Actual analytic
coordinate differentiation lowers original point-power order by at
most one, so $D_ig\in\mathfrak m^k$. Thus $a\in\mathfrak m^k$ as
well, because $b\in\mathfrak m^{k+1}\subset\mathfrak m^k$.

The actual Taylor derivative theorem identifies $T(D_ig)$ with the
literal formal partial derivative of $T(g)$. The preceding component
identity and true degree-$q$ agreement therefore identify
$\partial_i\operatorname{in}(F)$ with
$\operatorname{HC}_k(T(D_ig))$. The remainder $b$ has no degree-$k$
coefficient, so this equals $\operatorname{HC}_k(T(a))$. Since
$a\in I\cap\mathfrak m^k$, the preceding whole initial-component
membership theorem puts this exact component in the same $J$.
This proves derivative membership for \emph{every} original initial
generator, including all nonconvergent formal $F$ through actual
finite approximation.

Apply genuine ideal span induction. Additivity preserves derivative
membership; the actual Leibniz rule places the derivative of any
coefficient-times-generator expression in $J$ because both the
original generator and its derivative lie there. Hence the entire
same $J$ is stable under the actual partial derivative. Finally the
actual polynomial-series inclusion commutes with the literal
polynomial partial derivative, so apply this whole-ideal result to
the included polynomial. No selected initial family, proxy ideal or
word-coordinate basis is introduced.
\end{proof}

\begin{theorem}[Original HS-constant centre geometry yields genuine tangent polynomial partial stability at the same original origin ideal]
\label{hir:actual-original-hs-centre-whole-initial-tangent-pderiv}
\uses{hir:actual-original-first-symbol-whole-initial-pderiv-transfer,hir:actual-original-hs-centre-first-symbol-tangent-closure,hir:actual-original-plane-functional-power-germ-membership}
\lean{Hironaka.HironakaActualHomogeneousInitialTangentDerivatives.tangentInclusion_coordinateVector,Hironaka.HironakaActualHomogeneousInitialTangentDerivatives.exists_original_hs_constant_centre_initial_tangent_pderiv}
\leanok
From only the original coherent ideal, actual regular centre and
original intrinsic HS constancy, choose one genuine original
flattening and original generating family/opens before all orders,
as in the actual first-symbol theorem. For the literal same original
origin ideal $I_0$, its actual whole Taylor extension $K_0$ and the
existing whole homogeneous initial ideal $J_0$, every tangent
coordinate partial derivative preserves $J_0$. In particular the
literal polynomial $\operatorname{pderiv}_{i}(P)$ belongs to this
same initial ideal after genuine series inclusion whenever $P$ does,
for every actual tangent coordinate $i$.
This is an origin statement, not a common-neighborhood all-order
stability or normal-flatness theorem.
\end{theorem}
\begin{proof}
Take the genuine original all-positive-order first-symbol producer.
Its original flattening, original family/opens and origin ideal are
fixed before the universal order quantifier. At each positive order,
the actual tangent domain contains zero. Evaluate that order's
membership theorem at zero; the literal whole original translated-
sheaf ideal dictionary retains exactly $I_0$ in both input and output.
Thus the original first-symbol criterion holds at this same origin
for every positive order and every original tangent direction.

The actual tangent inclusion sends an original tangent coordinate
unit vector to that very ambient tangent coordinate vector: verify
its tangent and normal coordinates by the genuine inclusion formulas.
The canonical coordinate derivative in the original Taylor theorem
is therefore precisely the derivative appearing in this origin
criterion. Transport the original Taylor homomorphism and its genuine
coordinate directions together along the literal dimension equality
$0+n=n$. This is the accepted \texttt{actualUnshiftedTaylorAlgHom},
with the same original ideal retained as its input; equality induction
identifies each transported normal direction with the actual ambient
coordinate unit vector. It changes no ideal or Taylor coefficients. Apply the preceding whole-initial-ideal transfer for
each actual tangent coordinate. Its final polynomial conclusion is
already the literal polynomial partial derivative under the genuine
series inclusion, proving the asserted result. The per-order domains
have only been used at their common point zero; no common all-order
open, cylindrical normal-extension identity, normal-flatness or
compatible vertex-family theorem is inferred.
\end{proof}
